\documentstyle[10pt,amssymbols]{siamltex}
\newcommand{\dis}{\displaystyle}
\newcommand{\be}{\begin{equation}}
\newcommand{\ee}{\end {equation}}
%\renewcommand{\theequation}{{\arabic{equation}}}
\newcommand{\R}{\rm I\kern-.19emR}
\newcommand{\M}{\rm I\kern-.19emM}
\newcommand{\N}{\makebox[.80em]{\rm N\kern-.70emN}}
\newtheorem{example}[theorem]{Example}{\rm}
%\newtheorem{proposition}[theorem]{Proposition}{\rm}
\newcommand{\arw}{\mbox{$\rightarrow$}}
\newcommand{\vs}{\vspace{.15in}}
\newcommand{\vb}{\vspace{.4in}}
\newcommand{\vbb}{\vspace{1in}}
\newcommand{\mc}{\multicolumn}


\title{Gram polynomials and the Kummer function}


\author{R.W. Barnard\thanks{Department of Mathematics, Texas Tech
University, Lubbock, TX 79409}
\and
G. Dahlquist\thanks{Department of Computing Science, Royal Institute
of Technology, 10044 Stockholm, Sweden}
\and
K. Pearce\thanks{Department of Mathematics, Texas Tech University,
Lubbock, TX 79409}
\and
L. Reichel\thanks{Department of Mathematics and Computer Science,
Kent State University, Kent, OH 44242. Research supported in part by
NSF grants DMS-9205531 and DMS-9404706}
\and
K.C. Richards\thanks{Department of Mathematics, Southwestern
University, Georgetown, TX 78626}
}

\begin{document}

\maketitle

\begin{abstract}
Let $\left\{\phi_{k}\right\}^{n}_{k=0}$, $n<m$, be a family of polynomials
orthogonal with respect to the positive semi-definite bilinear form
$$
(g,h)_{d}:=\frac{1}{m}\sum^{m}_{j=1}g(x_{j})h(x_{j}),\qquad
x_{j}:=-1+(2j-1)/m.
$$
These polynomials are known as Gram polynomials. The present paper
investigates the growth of $|\phi_{k}(x)|$ as a function of $k$ and $m$
for fixed $x\in [-1,1]$. We show that when $n\leq 2.5m^{1/2}$, the
polynomials in the family $\left\{\phi_{k} \right\}^{n}_{k=0}$ are of
modest size on $[-1,1]$, and they are therefore well suited for the
approximation of functions on this interval. We also demonstrate
that if the degree $k$ is close to $m$, and $m \geq 10$, then $\phi_{k}(x)$
oscillates with large amplitude for values of $x$ near the endpoints of
$[-1,1]$, and this behavior makes $\phi_k$ poorly suited for the
approximation of functions on $[-1,1]$. We study the growth properties
of $|\phi_{k}(x)|$ by deriving a second order differential equation,
one solution of which exposes the growth. The
connection between Gram polynomials and this solution to the differential
equation suggested what became a long-standing conjectured inequality for
the confluent hypergeometric function ${}_1F_1$, also known as Kummer's
function, i.e., that
${}_1F_1((1-a)/2,1,t^2)\leq {}_1F_1(1/2,1,t^2)$ for all $a\geq 0$.
In this paper we completely resolve this conjecture by verifying a
generalization of the conjectured inequality with sharp constants.
\end{abstract}

\keywords
Confluent hypergeometric function, polynomial approximation
\endkeywords

%\AMS
%  xxxxxx
%\endAMS


\section{Introduction}

Let $f$ be a smooth function defined on the closed interval $[-1,1]$
and assume that $f$ is explicitly known only at the $m$ equidistant
points
\be\label{nodes}
x_{k} := -1 + (2k-1)/m,\qquad 1 \leq k \leq m.
\ee
We wish to approximate $f$ on $[-1,1]$ by a polynomial of degree $n$,
where $n < m$. Introduce the positive semi-definite bilinear form
\be \label{1.2}
(g,h)_{d} := \frac{1}{m} \dis\sum^{m}_{k=1} g(x_{k}) h(x_{k})
\ee
for functions $f,g$ continuous on $[-1,1]$, and define the associated discrete
semi-norm
\be \label{1.1}
\|g\|_{d} := (g,g)^{\frac{1}{2}}_{d}.
\ee
Let $\left\{\phi_{k} \right\}^{m-1}_{k=0}$ be the family of polynomials
that are orthogonal with respect to the bilinear form (\ref{1.2}), have
positive leading coefficient and are normalized so that
$\|\phi_{k}\|_d=1$. The $\phi_{k}$ are known as {\em Gram polynomials}.
These polynomials are discussed, e.g., by Dahlquist and Bj\"{o}rck
\cite[Section 4.4.4]{5}, Hildebrand \cite[Sections 7.13 and 7.16]{Hi} and
Szeg\"{o} \cite[Section 2.8]{14}.

Let $\Pi_n$ denote the set of all polynomials of degree at most $n$. The
polynomial $\Phi_n\in\Pi_n$, that solves the discrete least-squares
approximation problem
\be\label{least2}
\|f-\Phi_n\|_d=\min_{\Phi\in\Pi_n}\|f-\Phi\|_d,
\ee
is given by
\be \label{1.4}
\Phi_{n} (x) := \dis\sum^{n}_{k=0} \beta_{k} \phi_{k} (x),
\qquad \beta_{k} := ( \phi_{k}, f)_{d}
\ee
and is therefore simple to compute. It is the purpose of the present
paper to investigate the conditions on $n$ under which the solution
$\Phi_{n}$ of (\ref{least2}) also approximates $f$ well with respect
to the uniform norm
\[
\|g\|_{\infty} := \dis\sup_{x \in [-1,1]} |g(x)|.
\]
In order to gain some insight into the behavior of
$\|f-\Phi_{n}\|_{\infty}$,
we first review two special cases: $n \ll m$ and $n=m-1$. We begin with the
former. Let $\{p_{k}\}^{n}_{k=0}$ denote the {\em Legendre polynomials}
normalized so that $\|p_{k}\| =1$, where we define
\begin{eqnarray} \label{1.5}
<g,h> &:=& \frac{1}{2} \dis\int^{1}_{-1} g(x)h(x)dx, \\
\|g\| &:=& <g,g>^{1/2}
\end{eqnarray}
for all square integrable functions on $[-1,1]$. Analogously to (\ref{1.4}),
the solution $P_n\in\Pi_n$ of the (continuous) least-squares problem
\[
\|f-P_n\|=\min_{P\in\Pi_n}\|f-P\|
\]
can be written as
\[
P_{n} (x) = \dis\sum^{n}_{k=0} <p_{k}, f>p_{k}(x).
\]
In \cite[p. 345]{4} Brass proved the following result.

\begin{theorem}
Let $d\sigma$ be a distribution on $[-1,1]$, and let $\{q_k\}_{k=0}^{n+1}$
be a family of orthogonal polynomials with respect to $d\sigma$. Assume the
normalization $\int_{-1}^1 q_k^2(x) d\sigma(x)=1$. Let $d\sigma$ be such that\\

(i) $\int_{-1}^1 f(x) d\sigma(x) = \int_{-1}^1 f(-x) d\sigma(x)$ for any
$f\in C [-1,1]$,\\

(ii) $\|q_k\|_\infty=q_k(1)$, \quad $k=0,1,\dots,n+1$.\\

\noindent Assume that $f\in C^{n+1} [-1,1]$, and let
$\eta_k:=\int_{-1}^1 f(x)q_k(x)d\sigma(x)$. Then
$$
\|f-\sum^{n}_{k=0} \eta_k q_k \|_\infty\leq
\frac{\|q_{n+1}\|_\infty}{\|q^{(n+1)}_{n+1}\|_\infty}
\|f^{(n+1)}\|_\infty.
$$
Sharpness follows by letting $f=q_{n+1}$.
\end{theorem}

We apply this result and use the known properties of the Legendre
polynomials, including the fact that $\|p_k\|_\infty = p_k (1)$,
to obtain,
\be \label{1.8}
\|f-P_{n}\|_{\infty} \leq \dis\frac{\|f^{(n+1)}\|_{\infty}}{(n+1)!}
\|p_{n+1}\|_{\infty} \dis\lim_{x \rightarrow \infty} (x^{n+1}/p_{n+1}(x)).
\ee
This inequality is used in the proof of the following bound.

\begin{proposition}
Assume that $m>n$, and let $\Phi_n$ be given by (\ref{1.4}). Then, for
$f\in C^{n+1} [-1,1]$,
\be \label{1.11}
\|f-\Phi_{n}\|_{\infty} \leq \dis\frac{\|f^{(n+1)}\|_{\infty}}{2^{n}(n+1)!}
\cdot \dis\frac{\pi^{\frac{1}{2}}}{2} n^{\frac{1}{2}} (1+O(n^{-1})) +
\hat{c}_{n}O(m^{-2}),
\ee
where the $O(n^{-1})$-term is independent of $m$ and the $O(m^{-2})$-term
is independent of $n$. The constant $\hat{c}_{n}$ is independent of $f$ and
$m$.
\end{proposition}

\begin{proof}
The bilinear form (\ref{1.2}) corresponds to a discretization by the
rectangle rule of (\ref{1.5}), which has a discretization error
$O(m^{-2})$. Therefore, there are constants $c_{k}$, such that for each $k$,
\be \label{1.6}
\phi_{k}(x)=p_{k}(x) + c_{k} O(m^{-2}), \quad m \rightarrow \infty,
\ee
uniformly for $x \in [-1,1]$; see Wilson \cite{15} for details. It follows from
(\ref{1.6}) that there are constants $\hat{c}_{n}$, such that for each $n$,
\be \label{1.7}
\|f-\Phi_{n}\|_{\infty} \leq \|f-P_{n}\|_{\infty}+\|P_{n}-\Phi_{n}\|_{\infty}=
\|f-P_{n}\|_{\infty}+ \hat{c}_{n}O(m^{-2}), \qquad m \rightarrow \infty.
\ee
Substitute (\ref{1.8}) into (\ref{1.7}) and use the following
equalities that follow from results in \cite[Section 4.7]{14},
\begin{eqnarray}
\label{pnorm}
\|p_{n+1}\|_{\infty} &=& (2n+3)^{\frac{1}{2}}, \\
\label{1.10}
\dis\lim_{x \rightarrow \infty} (x^{n+1}/p_{n+1}(x))&=&2^{n}
{2n+1 \choose n}^{-1} (2n+3)^{-\frac{1}{2}},
\end{eqnarray}
and apply Stirling's formula to bound the binomial coefficient in (\ref{1.10}).
This shows the proposition.
\end{proof}

Let $Q_{n}\in\Pi_n$ solve the uniform-norm approximation problem
\[
\|f-Q_{n}\|_{\infty} = \min_{Q\in\Pi_n} \|f-Q\|_{\infty}.
\]
Then, for $f\in C^{n+1} [-1,1]$,
\be \label{1.12}
\|f-Q_{n}\|_{\infty} \leq \dis\frac{\|f^{(n+1)}\|_{\infty}}{2^{n}(n+1)!},
\ee
see Meinardus \cite[Theorem 60]{10}. The bound (\ref{1.12}) is sharp.
The closeness of the bounds (\ref{1.11}) and (\ref{1.12}) for large $m$,
suggests that for $m$ sufficiently large the polynomial $\Phi_n$,
given by (\ref{1.4}), is a good approximation of $f$ also when the
error is measured in the uniform norm.

We turn to the case when $n=m-1$. Then $\Phi_n$ interpolates $f$ at
the nodes (\ref{nodes}). A well-known difficulty arises: even for a
function $f$ analytic on $[-1,1]$, the approximant $\Phi_n$ may oscillate
with large amplitude near the endpoints of $[-1,1]$, and the amplitude may
increase with $n$. An analysis of this behavior, known as the {\em
Runge phenomenon}, is presented by Runge \cite{13}, and more recently
by Rivlin \cite{12} and Li and Saff \cite{LS}. The difficulty is
caused by the exponential growth with $n$ of the norm of the
interpolation operator; see \cite[p.\ 99]{12}.

A bound analogous to (\ref{1.8}) for Gram polynomials is shown to be
valid in Section 2. This suggests that $\Phi_n$, given by
(\ref{1.4}), approximates analytic functions $f$ well on $[-1,1]$ if the
degree $n$ is small enough in relation to $m$, so that
$\|\phi_{n+1}\|_{\infty}$ stays bounded as $n$ and $m$ increase. We
therefore need to study the growth of $\phi_{n}(x)$ as a function of
$m$, $n$ and $x$. In Section 3 we derive a family of second order
ordinary differential equations from the three-term recurrence
relation for the $\phi_{n}$. For each fixed value of $x\in [-1,1]$, we
obtain a differential equation that describes the behavior of
$\phi_{n}(x)$ for large values of $m$ and $n$. The differential
equation as well as the initial conditions on the solution depend on
the parameter $x\in [-1,1]$. The solution of each initial value problem
can be expressed in terms of the confluent hypergeometric function
\be\label{conhyp}
F(a,c,z):={ }_{1}F_{1}(a;c;z):=\sum_{n=0}^{\infty}\frac{(a)_{n}}
{(c)_{n}n!}z^{n},
\ee
where $(a)_n:=\Gamma(a+n)/\Gamma(a)$ is the Pochhammer symbol and
$\Gamma$ denotes the $\Gamma$-function. Different values of $x$
correspond to different values of the parameter $a$. The function
(\ref{conhyp}) is also known as Kummer's function.

Section 4 shows that the solution of the initial value problem corresponding
to $x=1$ dominates the solutions corresponding to $-1\leq x <1$. Therefore,
it suffices to consider only the former solution when studying the growth
of $\|\phi_n\|_{\infty}$ as $m$ and $n$ increase. The fact that the
solution corresponding to $x=1$ dominates solutions associated with the
other values of $x$ is equivalent to the inequality
\be\label{ineq}
F\left(\frac{1-\zeta}{2},1,z\right)\leq F(1/2,1,z)\;\;
\mbox{for all}\;\; \zeta\geq 0\;\; \mbox{and}\;\; z\geq 0.
\ee
The proof of (16) given in Section 4 is believed to be new.

Our study of solutions to the differential equation shows that for
large values of $m$ and $n$, the norm $\|\phi_{n}\|_{\infty}$ is
nearly invariant under changes in $n$ and $m$, whenever the ratio
$n/m^{\frac{1}{2}}$ is kept constant. Moreover, $\Phi_{n}$ defined by
(\ref{1.4}) is a good approximant of $f$ in the uniform norm,
provided that $n$ is not larger than a small multiple of
$m^{\frac{1}{2}}$, say $n \leq 2.5 m^{\frac{1}{2}}$.  Numerical
examples that illustrate the behavior of the Gram polynomials are
presented in Section 5.

The relevance of the ratio $n/m^{\frac{1}{2}}$ has previously been
noted by Bj\"{o}rk \cite{3} and Zaremba \cite{16} in their
investigation of Gram polynomials. Closely related problems are also
considered in \cite{6,7,8,Re,12}. Our method of investigation also can be used
to analyze classes of orthogonal polynomials other than Gram polynomials.


\section{Gram Polynomials}

The Gram polynomials introduced in Section 1 satisfy the
three-term recurrence relation, for $1\leq n < m$,
\begin{eqnarray} \label{2.1}
\phi_{n}(x) &=& 2 \alpha_{n-1} x \phi_{n-1}(x)-
\dis\frac{\alpha_{n-1}}{\alpha_{n-2}}\phi_{n-2} (x),  \\
\label{2.2}
\alpha_{n-1} &:=& \dis\frac{m}{n} \left(
\dis\frac{n^{2}- \frac{1}{4}}{m^{2}-n^{2}} \right)^{\frac{1}{2}},
\end{eqnarray}
with $\phi_{0}(x) := 1, \phi_{-1}(x) := 0$ and
$\alpha_{-1} := 1$; see, e.g., \cite[(4.4.24)-(4.4.26)]{5}.

\begin{theorem} \label{th1}
Each Gram polynomial $\phi_{n}$, $0 \leq n < m$, can be written as
a non-negative linear combination of Legendre polynomials $p_{j}$,
$0 \leq j \leq n$.  In particular,
\be \label{2.3}
\|\phi_{n} \|_{\infty} = \phi_{n}(1), \qquad 0 \leq n < m.
\ee
\end{theorem}

\begin{proof}
The theorem can be shown directly by induction.  It follows also from
a more general result by Askey \cite[Theorem 1]{2}. Here we verify
that the conditions of Theorem 1 in \cite{2} are satisfied.  Let
$\left\{ \phi^{\ast}_{n} \right\}^{m-1}_{n=0}$ be monic Gram
polynomials associated with the bilinear form (\ref{1.2}), and let
$\left\{p^{\ast}_{n} \right\}^{\infty}_{n=0}$ be monic Legendre
polynomials.  Then
\begin{eqnarray*}
\phi^{\ast}_{n+1}(x) &=& x \phi^{\ast}_{n}(x)-\lambda_{n}\phi^{\ast}_{n-1}(x),
\qquad 0 \leq n \leq m-2, \\
p^{\ast}_{n+1}(x) &=& x p^{\ast}_{n}(x) - \delta_{n} p^{\ast}_{n-1} (x),
\qquad n=0,1,\ldots~,
\end{eqnarray*}
where
\begin{eqnarray*}
\lambda_{n} &:=& \dis\frac{n^{2}}{4n^{2}-1}
( 1- \dis\frac{n^{2}}{m^{2}} ), \\
\delta_{n} &:=& \dis\frac{n^{2}}{4n^{2}-1}
\end{eqnarray*}
and $\phi^*_{0} (x) := p^{\ast}_{0} (x) := 1$,
$\phi^*_{-1} (x):=p^{\ast}_{-1}(x) := 0$. We have to show that
$\delta_k\geq\lambda_n>0$ for $1\leq k\leq n$ and $0\leq n\leq m-2$. But
$\delta_k$ decreases as $k\geq 1$ increases, and $\delta_n\geq\lambda_n>0$.
Thus, the conditions of \cite[Theorem 1]{2} are satisfied, and, therefore,
\be\label{expan}
\phi^{\ast}_{n} (x) = \dis\sum^{n}_{j=0} a_{nj} p^{\ast}_{j}(x)
\ee
with $a_{nj} \geq 0$ for all $0 \leq j \leq n$ and $0\leq n <m$. The inequality
(\ref{2.3}) now follows from the representation (\ref{expan}) and the
fact that $\|p^*_j\|_\infty = p^*_j(1)$.
\end{proof}

It follows from (\ref{2.3}) and Theorem 1.1 that the error bound
\be\label{bound}
\|f- \Phi_{n}\|_{\infty} \leq \dis\frac{\|f^{(n+1)}\|_{\infty}}{(n+1)!}
\left\| \phi_{n+1} \right\|_{\infty} \dis\lim_{x \rightarrow \infty} (x^{n+1}/
\phi_{n+1} (x)),
\ee
which is analogous to (\ref{1.8}), is valid. Substitution of
$f := \phi_{n+1}$ into (\ref{bound}) shows the sharpness of the bound.


\section{A Differential Equation Model}

A differential equation is derived that approximates the three-term
recurrence relation for $\phi_{n}(x)$.  The solution of the differential
equation is a function of $(n- \frac{1}{2})/m^{\frac{1}{2}}$. In order
to derive the differential equation, we first introduce
$\tau :=n/m^{\frac{1}{2}}$. Then (\ref{2.2}) can be written as
\be \label{3.1a}
\alpha_{n-1} = (1- \dis\frac{1}{4} \tau^{-2} m^{-1})^{\frac{1}{2}}
(1- \tau^{2} m^{-1})^{-\frac{1}{2}}.
\ee
Let $\tau_0$ and $\tau_1$ be constants, such that
$0 < \tau_{0} < \tau_1 < \infty$, and consider $n$ as a function of $\tau$.
We obtain from (\ref{3.1a}) that
\be\label{3.1}
\alpha_{n-1} = 1 + \dis\frac{1}{2}(\tau^{2}- \dis\frac{1}{4}\tau^{-2}) m^{-1}
+ O(m^{-2}), \qquad m \rightarrow \infty,
\ee
where the convergence in (\ref{3.1}) is uniform for $\tau_0\leq\tau\leq\tau_1$.
Note that the bound $n\geq \tau_0 m^{1/2}$ implies that $n\rightarrow\infty$
as $m\rightarrow\infty$. From (\ref{3.1a}), we also obtain
\be \label{3.2}
\dis\frac{\alpha_{n-1}}{\alpha_{n-2}} =1+ ( \tau + \dis\frac{1}{4}
\tau^{-3} ) m^{-\frac{3}{2}} + O (m^{-2}),\qquad m \rightarrow \infty,
\ee
uniformly for $\tau_{0} \leq \tau \leq \tau_{1}$. Let $x := 1-\zeta/m$. Then
(\ref{2.1}) can be written in the form
\begin{eqnarray}
\nonumber
 & &\dis\frac{\phi_{n}(x)-2 \phi_{n-1}(x) + \phi_{n-2}(x)}{m^{-1}} \\
\label{3.3}
 &-& 2m( \alpha_{n-1}-1)\phi_{n-1}(x)
     -m ( 1- \dis\frac{\alpha_{n-1}}{\alpha_{n-2}} ) \phi_{n-2} (x) \\
\nonumber
 &+& 2 \zeta \alpha_{n-1} \phi_{n-1} (x) =0.
\end{eqnarray}
Substituting (\ref{3.1}) and (\ref{3.2}) into (\ref{3.3}) yields
\begin{eqnarray}
\nonumber
 & &\dis\frac{\phi_{n}(x)-2 \phi_{n-1}(x) + \phi_{n-2}(x)}{m^{-1}} \\
 \label{3.4}
 &=& (\tau^{2}- \frac{1}{4} \tau^{-2} -2 \zeta) \phi_{n-1} (x) -
    ( \tau + \dis\frac{1}{4} \tau^{-3}) m^{-\frac{1}{2}} \phi_{n-2}(x) \\
\nonumber
 &+& \phi_{n-1}(x) O(m^{-1}) + \phi_{n-2}(x) O(m^{-1}),
\qquad m \rightarrow \infty,
\end{eqnarray}
where the convergence is uniform for $\tau_{0} \leq \tau \leq \tau_{1}$.
Introduce $t := \tau - \frac{1}{2} \Delta \tau$,
$\Delta \tau := \Delta t := m^{-\frac{1}{2}}$ and substitute
\be \label{3.5}
\phi (t) := \phi_{n-1}(x)/ \sqrt{2m^{\frac{1}{2}}}
\ee
into (\ref{3.4}).  The change of variables from
$\tau$ to $t$ makes the $O(m^{-\frac{1}{2}})$-term vanish.  We
obtain
\[
\dis\frac{\phi(t+ \Delta t)-2 \phi(t) +\phi(t- \Delta t)}{(\Delta t)^{2}}
= ( t^{2} - \dis\frac{1}{4}t^{-2} - 2 \zeta)
\phi (t) + O (\Delta t^{2}), \qquad \Delta t \rightarrow 0,
\]
and, hence,
\be \label{3.6}
\frac{d^2}{dt^2}\phi(t) = (t^{2} - \dis\frac{1}{4} t^{-2} -2 \zeta)
\phi(t) + O (\Delta t^{2}), \qquad \Delta t \rightarrow 0,
\ee
The convergence in (\ref{3.6}) is
uniform for $t_{0} \leq t \leq t_{1}$, where $t_{0}, t_{1}$ are
arbitrary but fixed constants, such that $0 < t_{0} < t_1 < \infty$.  From
(\ref{3.6}) we obtain the differential equation
\be \label{3.7}
\frac{d^2}{dt^2}\phi(t) = (t^{2} - \dis\frac{1}{4} t^{-2} -2 \zeta) \phi (t).
\ee
The general solution of (\ref{3.7}) is given by
\be \label{3.8}
\phi(t) = t^{\frac{1}{2}} e^{-t^{2}/2}
(A\,\, {}_1F_1 ( \frac{1}{2} (1- \zeta), 1, t^{2} ) +
B\, U( \frac{1}{2}(1 - \zeta ), 1, t^{2})),
\ee
where $A,B$ are arbitrary constants, $F={}_1F_1$ is Kummer's function
(\ref{conhyp}), and $U$ is a linearly independent logarithmic solution to
Kummer's equation; see \cite[p. 504]{1} for the definition of $U$. The
differential equation model (\ref{3.6}), the solution (\ref{3.8}) with
$A=1$ and $B=0$, and equation (\ref{lagu}) below were first suggested in
\cite{DR}.

We are interested in studying $\|\phi_{n}\|_{\infty} = \phi_{n}(1)$,
and therefore choose $\zeta = 0$ in (\ref{3.7}) and (\ref{3.8}). This
value of $\zeta$ corresponds to $x=1$. Other choices
of $\zeta$ are discussed below. For $\zeta = 0$, the solution (\ref{3.8})
simplifies to, see \cite [(13.6)]{1},
\[
\phi(t) = t^{\frac{1}{2}} (A I_{0}(t^{2}/2) + B \pi^{-\frac{1}{2}} K_{0}
(t^{2}/2)),
\]
where $I_{0}$ and $K_{0}$ are modified Bessel functions of zeroth
order of the first and second kind, respectively.  We note that, see
\cite[Chapter 9]{1},



\[
\begin{array}{rl}
t^{\frac{1}{2}}I_{0}(t^{2}/2) = \pi^{-\frac{1}{2}} t^{-\frac{1}{2}}
e^{t^{2}/2} (1 + O (t^{-2})), &  t \rightarrow \infty,\\
t^{\frac{1}{2}} K_{0} (t^{2}/2)= \pi^{\frac{1}{2}} t^{-\frac{1}{2}}
e^{-t^{2}/2} (1 + O (t^{-2})), & t \rightarrow \infty,
\end{array}
\]
which shows that $t^{\frac{1}{2}} I_{0} (t^{2}/2)$ is a dominating solution
of (\ref{3.7}) as $t$ increases. Moreover,
\be \label{3.10}
\begin{array}{ll}
t^{\frac{1}{2}}I_{0}(t^{2}/2) = t^{\frac{1}{2}}
(1 + O (t^{4})), &  t \rightarrow 0,\\
t^{\frac{1}{2}} K_{0} (t^{2}/2)= t^{\frac{1}{2}} ((-2 \ln (t/2) +
\gamma) I_{0} (t) + O(t^{4}) ), & t \rightarrow 0,
\end{array}
\ee
where $\gamma \approx 0.577$ denotes Euler's constant.

We turn to the initial conditions.  Since $\|p_{n-1}\|=p_{n-1}(1)$,
we obtain from (\ref{1.6}) and (\ref{pnorm}), that for fixed $n$,
\be \label{3.11}
\phi_{n-1}(1) = (2n+1)^{\frac{1}{2}} + c_{n-1} O(m^{-2}) =
\sqrt{2m^{\frac{1}{2}}} t^{\frac{1}{2}} +c_{n-1} O(m^{-2}),
\qquad m \rightarrow \infty.
\ee
Substituting (\ref{3.11}) into (\ref{3.5}) yields
\[
\phi(t) = t^{\frac{1}{2}} (1 + O (t^{4})),\qquad t \rightarrow 0,
\]
and in view of (\ref{3.10}), we obtain
\be \label{3.12}
\phi(t) - t^{\frac{1}{2}} I_{0} (t^{2}/2) =
t^{\frac{1}{2}} O (t^{4}), \qquad t \rightarrow 0,
\ee
where the power of $t$ in the $O(t^{4})$-factor cannot be increased. Thus,
the function
\be\label{phi0}
\phi^{(0)}(t):=t^{\frac{1}{2}} I_{0}(t^{2}/2)
\ee
can be used to approximate $\phi_{n}(1)/\sqrt{2m^{\frac{1}{2}}}$ in the
following way. Let $\phi(t)$ be defined by (\ref{3.5}) with $x=1$, and
select $t_{0} > 0$ sufficiently small so that the right-hand side of
(\ref{3.12}) is small for $0 \leq t \leq t_{0}$. Analogously to (\ref{3.5}),
define
\[
\hat{\phi}_{n-1}(x):=\sqrt{2m^\frac{1}{2}} \phi^{(0)}(t),\qquad
t:=(n-\dis{\frac{1}{2}})/m^\frac{1}{2},\qquad t_0\leq t\leq t_1,
\]
and choose $m$ large enough so that $\hat{\phi}_{n-1}(1)$ is a good approximate
solution of the difference equation (\ref{3.3}) for $t_{0} \leq t \leq t_{1}$.
Since $\phi^{(0)}$ is a dominating solution of (\ref{3.7}),
it models the behavior of the scaled polynomials $\phi_{n}(x)/
\sqrt{2m^{\frac{1}{2}}}$ at $x=1$ fairly well already for modest
values of $m$. This is illustrated by numerical examples in Section 5.

We next determine initial conditions for $\zeta > 0$.  For bounded
$\zeta > 0$ and fixed $n$, we obtain, by (\ref{1.6}), that
\begin{eqnarray*}
\phi_{n-1}(1- \zeta/m) & = & p_{n-1}(1-\zeta/m) + c_{n-1} O (m^{-2})  \\
 & = &  p_{n-1}(1) + \tilde{c}_{n-1} O (m^{-1}), \qquad m \rightarrow \infty,
\end{eqnarray*}
where the constant $\tilde{c}_{n-1}$ is independent of $m$. For $\phi(t)$
defined by (\ref{3.5}), with $x=1-\zeta/m$, we have
$\phi(t) = t^{\frac{1}{2}} (1 + O(t^{2}))$, $t \rightarrow 0$.
Analogously to (\ref{3.12}), we find that
\[
\phi(t) -t^{\frac{1}{2}} e^{-t^{2}/2}F
( \dis\frac{1}{2} (1-\zeta), 1, t^{2}) = t^{\frac{1}{2}} O(t^{2}),
\qquad t \rightarrow 0.
\]
The solution of (\ref{3.7}) that models the behavior of
$\phi_{n-1} (1-\zeta/m) / \sqrt{2m^{\frac{1}{2}}}$ for $\zeta>0$ is therefore
\be \label{phizeta}
\phi^{(\zeta)} (t) := t^{\frac{1}{2}}
e^{-t^{2}/2} F( \dis\frac{1}{2} (1-\zeta), 1, t^{2}).
\ee
Note that $\phi^{(\zeta)}(t)\rightarrow\phi^{(0)}(t)$ as $\zeta\rightarrow 0$.
The fact that $\|\phi_n\|_\infty=\phi_n(1)$ suggests the inequality
\be\label{polineq}
\phi^{(\zeta)}(t)\leq \phi^{(0)}(t)\ \ \mbox{for~all}\ \ \zeta\geq 0,\ \
t\geq 0,
\ee
which is equivalent with (\ref{ineq}). We will show (\ref{polineq}) in
Section 4.

Let $\zeta := 2k+1$ for some integer $0\leq k < m$. Then
$x:=1-\zeta/m$ is the node $x_{m-k}$ defined by (\ref{nodes}), and we
obtain from (\ref{phizeta}) that the solution
\be\label{lagu}
\phi^{(\zeta)}(t) = t^{\frac{1}{2}} e^{-t^{2}/2} F(-k, 1, t^{2}) =
t^{\frac{1}{2}} e^{-t^{2}/2} L_{k} (t^{2})
\ee
of (\ref{3.7}) models the behavior of
$\phi_{n-1} (1-\zeta/m) / \sqrt{2m^{\frac{1}{2}}}$. Here $L_{k}(x)$ denotes
a Laguerre polynomial of degree $k$; see \cite[(22.5.54)]{1}. The
fact that $\phi^{(\zeta)}(t) \rightarrow 0$ as $t \rightarrow \infty$
agrees well with the observed behavior of the polynomials $\phi_{n-1}$
at the nodes; see Table 3 of Section 5.


\section{An Inequality for Kummer's function}

Inequality (\ref{2.3}), the connection between Gram polynomials and
the confluent hypergeometric function exposed in Section 3, and
numerical evidence suggested the stronger inequality (\ref{ineq}). The
latter inequality was first presented as a conjecture in 1985 \cite{Ba0},
and is also discussed in \cite[p.\ 21]{Ba}. For completeness, and because of
its independent interest, we verify a generalization of this conjecture with
sharp constants.

\begin{theorem}\label{theo1} For all $\zeta\geq 0$, $x\geq 0$, and
$c\geq 1/2$
\begin{equation}\label{4x2}
F\left(\frac{1-\zeta}{2},c,x\right)\leq F(1/2,c,x).
\end{equation}
Moreover, $c\geq 1/2$ is sharp, i.e., for $c<1/2$ and $x>0$, there is a
$\zeta>0$ such that inequality (\ref{4x2}) fails.
\end{theorem}

\begin{proof}
Suppose that $x>0$ and $\zeta>0$, and consider the special case $c=1/2$.
We will make use of the following classical identity, see
\cite[p. 1085, \#9.211-3]{GR},
\begin{equation}\label{35}
F(-\nu,\alpha+1,x)=\frac{\Gamma(\alpha+1)}{\Gamma(\alpha+\nu+1)}
e^x x^{-\frac{\alpha}{2}} \int_0^\infty e^{-t}
t^{\nu+\frac{\alpha}{2}} J_\alpha (2\sqrt{xt}) dt,
\end{equation}
for $\alpha+\nu+1>0$, where $J_\alpha$ is the Bessel function of order
$\alpha$.  Using identity (\ref{35}) with $\nu=\frac{\zeta-1}{2}$ and
$\alpha=-1/2$, it follows that $\alpha+\nu+1=\zeta/2>0$, and
$$
J_\alpha(z)=J_{-\frac{1}{2}}(z)=\sqrt{\frac{2}{\pi z}}\cos(z),\;\;
$$
see \cite[p.\ 977, \#8.464-2]{GR}. Therefore,
\begin{eqnarray*}
F\left(\frac{1-\zeta}{2},1/2,x\right) & = &
e^x\frac{\Gamma(1/2)}{\Gamma(\zeta/2)} x^{\frac{1}{4}}
\int_0^\infty e^{-t}t^{\frac{\zeta}{2}-\frac{3}{4}}
J_{-1/2}(2\sqrt{xt}) dt \\[1ex]
& = & e^x\frac{\Gamma(1/2)}{\Gamma(\zeta/2)} x^{\frac{1}{4}}
\int_0^\infty e^{-t}t^{\frac{\zeta}{2}-\frac{3}{4}}
\sqrt{\frac{2}{2\pi\sqrt{xt}}}\cos(2\sqrt{xt}) dt \\[1ex]
%% & = & e^x\frac{\sqrt{\pi}}{\Gamma(\zeta/2)}x^{\frac{1}{4}}
%% \int_0^\infty e^{-t}t^{\frac{\zeta}{2}-\frac{3}{4}}
%% \frac{\cos(2\sqrt{xt})}x^{\frac{1}{4}}t^{\frac{1}{4}}\sqrt{\pi}
%% dt \\[1ex]
& = & e^x\frac{1}{\Gamma(\zeta/2)}\int_0^\infty e^{-t}
t^{\frac{\zeta}{2}-1}\cos(2\sqrt{xt}) dt \\[1ex]
& \leq & e^x\frac{1}{\Gamma(\zeta/2)}\int_0^\infty e^{-t}
t^{\frac{\zeta}{2}-1} dt \\[1ex]
& = & e^x = F(1/2,1/2,x),
\end{eqnarray*}
which implies that
\begin{equation} \label{36}
F\left(\frac{1-\zeta}{2},c,x\right)\leq F(1/2,c,x)\;\;
\mbox{for}\;\; c=1/2\;\; \mbox{and for all}\;\; x\geq 0\;\;
\mbox{and}\;\; \zeta\geq 0.
\end{equation}

Now suppose that $c>1/2$. We wish to relate $F(a,c,x)$ to $F(a,1/2,x)$.
The identity, see \cite[p. 863, \#7.613-1]{GR},
\[
F(a,c,x)=\frac{\Gamma(c)x^{1-c}}{\Gamma(\gamma)\Gamma(c-\gamma)}
\int_0^x t^{\gamma-1}(x-t)^{c-\gamma-1} F(a,\gamma,t)\,dt\;\;
\mbox{for}\;\; c>\gamma>0,
\]
with $x>0$, $\zeta>0$ and $c>\gamma=1/2$, yields
\begin{eqnarray}
\nonumber
F\left(\frac{1-\zeta}{2},c,x\right)
& = & \frac{\Gamma(c)x^{1-c}}{\Gamma(1/2)\Gamma(c-1/2)}
\int_0^x t^{-\frac{1}{2}}(x-t)^{c-\frac{3}{2}}
F\left(\frac{1-\zeta}{2},1/2,t\right) dt \\[1ex]
\label{6}
& \leq & \frac{\Gamma(c)x^{1-c}}{\Gamma(1/2)\Gamma(c-1/2)}
\int_0^x t^{-\frac{1}{2}}(x-t)^{c-\frac{3}{2}}
F(1/2,1/2,t) dt \\[1ex]
\nonumber
& = & F(1/2,c,x),
\end{eqnarray}
where the inequality in (\ref{6}) follows from (\ref{36}) and the fact that
$$
\frac{\Gamma(c)x^{1-c}}{\Gamma(1/2)\Gamma(c-1/2)} t^{-\frac{1}{2}}
(x-t)^{c-\frac{3}{2}}>0
$$
for $x>t\geq 0$ and $c>1/2$. This establishes that
\begin{equation}\label{39}
F\left(\frac{1-\zeta}{2},c,x\right)\leq F(1/2,c,x)\;\; \mbox{for
all}\;\; \zeta\geq 0,x\geq 0,\;\; \mbox{and}\;\; c\geq 1/2.
\end{equation}
The sharpness of $c=1/2$ will follow from (the proof of) Theorem~\ref{theo2}.
\end{proof}

Another concise version of inequality (\ref{39}) is revealed when
it is expressed in terms of the Whittaker functions $M_{\lambda,\mu}$,
which are given by, see \cite[p.\ 1087, \#9.220-2]{GR},
\[
M_{\lambda,\mu}(x):=x^{\mu+\frac{1}{2}} e^{-x/2} F
\left(\mu-\lambda+\frac{1}{2},1+2\mu,x\right).
\]

\begin{theorem}\label{theo2} Suppose that $\lambda\geq\mu\geq-1/4$.
Then for all $x\geq 0$,
\begin{equation}\label{whit}
M_{\lambda,\mu}(x)\leq M_{\mu,\mu}(x).
\end{equation}
Moreover, $\mu=-1/4$ is sharp, i.e., for any $\mu<-1/4$ and $x>0$,
there is a $\lambda>\mu$ such that inequality (\ref{whit}) is invalid.
\end{theorem}

\begin{proof} Suppose that $-1/4\leq\mu<\lambda$ and $x>0$.  For
$c=1+2\mu$ and $\frac{1-\zeta}{2}=\mu-\lambda+\frac{1}{2}$ it follows
that $c\geq 1/2$ and $\zeta=2(\lambda-\mu)>0$. We have
\begin{eqnarray}
\nonumber
M_{\lambda,\mu}(x) & = & x^{\mu+\frac{1}{2}} e^{-x/2} F
\left(\mu-\lambda+\frac{1}{2},1+2\mu,x\right) \\[1ex]
\nonumber
& = & x^{\frac{c}{2}} e^{-x/2} F
\left(\frac{1-\zeta}{2},c,x\right) \\[1ex]
\label{10}
& \leq & x^{\frac{c}{2}} e^{-x/2} F (1/2,c,x)\\[1ex]
\nonumber
& = & M_{\mu,\mu}(x),
\end{eqnarray}
where the inequality (\ref{10}) follows from (\ref{39}). Therefore,
inequality (\ref{whit}) holds for all $\lambda\geq\mu\geq-1/4$ and $x\geq 0$.

In order to demonstrate the sharpness of $\mu=-1/4$, we note the
asymptotic relationship for large $\lambda>0$ given by, see
\cite[p. 1089, \#9.228]{GR},
\[
M_{\lambda,\mu}(x)\sim\frac{\Gamma(1+2\mu)}{\sqrt{\pi}}
\lambda^{-\mu-\frac{1}{4}}x^{\frac{1}{4}}
\cos(2\sqrt{\lambda x}-\mu\pi-\frac{\pi}{4}).
\]
Now let $x_0>0$ and $\mu_0\in(-1/2,-1/4)$ both be fixed.
For each positive integer $n$, let $\lambda_n$ satisfy
$$
2\sqrt{\lambda_nx_0}-\mu_0\pi-\pi/4=2n\pi,\;\; \mbox{i.e.,}\;\;
\lambda_n=(2n\pi+\mu_0\pi+\frac{\pi}{4})^2/(4x_0).
$$
Since $-\mu_0-\frac{1}{4}>0$, it follows that
$\lambda_n^{-\mu_0-\frac{1}{4}}\rightarrow \infty$ as
$n\rightarrow \infty$. Thus, the expression
$\frac{\Gamma(1+2\mu_0)}{\sqrt{\pi}}\lambda_n^{-\mu_0-\frac{1}{4}}
x_0^{\frac{1}{4}}$ can be made arbitrarily large by choosing a
sufficiently large positive integer $n$. In particular, there is a
$\lambda_n$ sufficiently large, such that
\begin{eqnarray*}
M_{\lambda_n,\mu_0}(x_0) & \sim & \frac{\Gamma(1+2\mu_0)}
{\sqrt{\pi}}\lambda_n^{-\mu_0-\frac{1}{4}}x_0^{\frac{1}{4}}\cos
\left(2\sqrt{\lambda_nx_0}-\mu_0\pi-\frac{\pi}{4}\right)\\[1ex]
& = & \frac{\Gamma(1+2\mu_0)}{\sqrt{\pi}}\lambda_n^{-\mu_0-\frac{1}
{4}}x_0^{\frac{1}{4}}\cos(2n\pi)\\[1ex]
& = & \frac{\Gamma(1+2\mu_0)}{\sqrt{\pi}}\lambda_n^{-\mu_0-\frac{1}
{4}}x_0^{\frac{1}{4}}\cdot 1\\[1ex]
& > & M_{\mu_0,\mu_0}(x_0).
\end{eqnarray*}

Therefore, for each $x>0$ and $\mu\in(-1/2,-1/4)$, there is a
$\lambda>\mu$, such that $M_{\lambda,\mu}(x)>M_{\mu,\mu}(x)$.  This
proves the sharpness of $\mu=-1/4$, and, hence, the sharpness of
$c=1+2\mu=1/2$ in Theorem~\ref{theo1}.
\end{proof}


\section{Numerical Examples}

The behavior of the Gram polynomials $\phi_{n}$ is displayed in three
tables. The tables compare $\phi_n(x)$, for several values of $x$, with
the function
$\phi^{(\zeta)}(t)$, which is given by either (\ref{phi0}), (\ref{phizeta}) or
(\ref{lagu}) depending on the value of $\zeta$. Throughout this section
$x:=1-\zeta/m$ and $t:=(n-\frac{1}{2})/m^\frac{1}{2}$. We use the
notation $M (E)$ for the number $M\cdot 10^E$ in the tables.
All computations were carried out on a VAX 11/780 computer in double
precision arithmetic, i.e., with about $15$ significant digits.

Table 1 shows the error
$(\phi_{n-1}(x)/\sqrt{2m^{\frac{1}{2}}} - \phi^{(\zeta)}(t))/\sqrt{t}$ for
$\zeta=0$ and $\zeta=1/2$. Columns 4 and 5 show that
$\phi_{n-1}(1)>\phi_{n-1}(1-\frac{1}{2}m)$, in agreement with our analysis.
Columns 6 and 7 illustrate the convergence of the error
$(\phi_{n-1}(x)/\sqrt{2m^{\frac{1}{2}}} - \phi(t))/\sqrt{t}$ as $m$
increases and $t$ is in a fixed interval. Note that the error is positive.

Table 2 displays the rapid growth of $\phi^{(0)}(t)$ with $t$. Recall that
$\sqrt{2m^{1/2}}\phi^{(0)}(t)$ approximates $\phi_n(1)$ for
$t=(n-\frac{1}{2})/m^{1/2}$. The fast growth of $\phi^{(0)}(t)$ with $t$
indicates that $\phi_n(1)$ grows rapidly with
$t=(n-\frac{1}{2})/m^{\frac{1}{2}}$.
The table suggests that the choice of $m$ and $n$ should be such that
$t=(n-\frac{1}{2})/m^{\frac{1}{2}} \leq 2.5$ in order to keep the norm
$\|\phi_n\|_\infty$ modest. The norm $\|\phi_{n-1}\|=\phi_{n-1}(1)$
can be determined from Table 1 for such values of $m$ and $n$.

Table 3 shows the behavior of $\phi_{n-1}(x_m)/\sqrt{2m^{\frac{1}{2}}}$,
where the node $x_m$ is defined by (\ref{nodes}). Thus, $\zeta=1$. The
table shows that both $\phi^{(1)} (t)$ and
$\phi_{n-1} (x_m) / \sqrt{2m^{\frac{1}{2}}}$ are small for large values of $t$.

\vskip10pt
\noindent
{\bf Acknowledgement}. L.R. would like to thank Bill Gragg for helpful
discussions.

\newpage
\begin{thebibliography}{18}
\bibitem{1}
{\sc M. Abramowitz and I. A. Stegun}, {\em Handbook of Mathematical
Functions}, NBS Applied Mathematics Series, 55, 1968.

\bibitem{2}
{\sc R. Askey}, {\em Orthogonal expansions with positive coefficients II},
SIAM J. Math. Anal., 2 (1971), 340-346.

\bibitem{Ba0}
{\sc R. W. Barnard}, {\em Survey of open problems and conjectures in
complex analysis and special functions}. Talk presented at the
Symposium on the Proof of the Bieberbach Conjecture, Purdue
University, West Lafayette, 1985.

\bibitem{Ba}
{\sc R. W. Barnard}, {\em Open problems and conjectures in complex
analysis}, Computational Methods and Function Theory, eds. S.
Ruscheweyh, E. B. Saff, L. C. Salinas, R. S. Varga, Lecture Notes in
Mathematics \# 1435, Springer, New York, 1990.

\bibitem{3}
{\sc H. Bj\"{o}rk}, {\em Contribution to the problem of least squares
approximation}, Report Na 71.37, Dept. of Numerical Analysis, Royal
Institute of Technology, Stockholm, 1971.

\bibitem{4}
{\sc H. Brass}, {\em Error estimates for least squares approximation
by polynomials}, J. Approx. Theory, 41 (1984), 345-349.

\bibitem{5}
{\sc G. Dahlquist and \AA. Bj\"{o}rck}, {\em Numerical Methods}, Prentice Hall,
Englewood Cliffs, 1974.

\bibitem{DR}
{\sc G. Dahlquist and L. Reichel}, {\em Polynomial approximation studied by
a differential equation model}, Unpublished lecture notes, Stockholm and
Moscow, Aug. 1984.

\bibitem{6}
{\sc H. Ehlich}, {\em Polynome zwischen Gitterpunkten}, Math. Z.,
93 (1966), 144-153.

\bibitem{7}
{\sc H. Ehlich and K. Zeller}, {\em Schwankung von Polynomen zwischen
Gitterpunkten}, Math. Z., 86 (1964), 41-44.

\bibitem{GR} {\sc I. S. Gradshteyn and I. M. Ryzhik}, {\em Tables of
Integrals, Series, and Products}, 5th ed., Academic Press, Orlando,
1994.

\bibitem{Hi} {\sc F. B. Hildebrand}, {\em Introduction to Numerical Analysis},
Dover, New York, 1987.

\bibitem{8}
{\sc A. Kro\'{o}}, {\em A comparison of uniform and discrete polynomial
approximation}, Analysis Mathematica, 5 (1979), 35-49.

\bibitem{LS}
{\sc X. Li and E. B. Saff}, {\em Local convergence of Lagrange interpolation
associated with equidistant nodes}, J. Approx. Theory, 78 (1994), 213-225.

\bibitem{10}
{\sc G. Meinardus}, {\em Approximation of Functions:  Theory and
Numerical Methods}, Springer, New York, 1967.

\bibitem{Re}
{\sc L. Reichel}, {\em On polynomial approximation in the uniform norm by
the discrete least squares method}, BIT, 26 (1986), 349-368.

\bibitem{11}
{\sc T. J. Rivlin}, {\em An Introduction to the Approximation of Functions},
Blaisdell, Waltham, 1966.

\bibitem{12}
{\sc T. J. Rivlin and E. W. Cheney}, {\em A comparison of uniform
approximation on an interval and a finite subset thereof}, SIAM J.
Numer. Anal., 3 (1966), 311-320.

\bibitem{13}
{\sc C. Runge}, {\em \"{U}ber empirische Funktionen und die Interpolation
zwischen \"{a}quidistanten Ordinaten}, Z. f. Math., 46 (1901), 224-243.

\bibitem{14}
{\sc G. Szeg\H{o}}, {\em Orthogonal Polynomials}, 4th ed., Amer. Math.
Soc., Providence, RI, 1975.

\bibitem{15}
{\sc M. W. Wilson}, {\em Convergence properties of discrete analogs
of orthogonal polynomials}, Computing 5 (1970), 1-5.

\bibitem{16}
{\sc S. K. Zaremba}, {\em Some properties of polynomials orthogonal over
the set $< 1, 2, \ldots, N>$}, Annali Matem., Ser. 4, 105
(1975), 333-345.
\end{thebibliography}

\newpage

\begin{table}
\caption{Accuracy for increasing $m$ for $t$ in a fixed interval; $x=1-\zeta/m$}
\begin{center}
\begin{tabular}{|ccc|cc|cc|}
\hline
& & & & & & \\
& & & \mc{2}{|c|}{$\phi_{n-1}(x)/\sqrt{2m^{\frac{1}{2}}}$} &
\mc{2}{c|}{$\left(\dis\frac{\phi_{n-1} (x)}{\sqrt{2m^{\frac{1}{2}}}}
- \phi^{(\zeta)} (t) \right) / \sqrt{t}$} \\
& & & & & & \\
$m$ & $n-1$ & $t$ & $\zeta=0$ & $\zeta=\frac{1}{2}$ & $\zeta=0$ &
$\zeta=\frac{1}{2}$ \\
& & & & & & \\
\hline
$20$ & $\phantom{1}1$  & $0.34$ & $5.80 (-1)$ & $5.65 (-1)$ & $4.61 (-4)$ &
$3.37 (-3)$\\
$20$ & $\phantom{1}5$ & $1.23$ & $1.29\phantom{(-1)}$ & $8.67 (-1)$ &
$1.22 (-2)$ & $6.53 (-3)$ \\
$20$ & $10$ & $2.35$ & $7.08\phantom{(-1)}$ &$2.48\phantom{(-1)}$ &
$6.05 (-1)$ & $1.99 (-1)$ \\
\hline
$40$ & $\phantom{1}1$  & $0.24$ & $4.87 (-1)$ & $4.81 (-1)$ & $1.15 (-4)$ &
$1.62 (-3)$ \\
$40$ & $\phantom{1}5$ & $0.87$ & $9.68 (-1)$ & $7.97 (-1)$ & $2.06 (-3)$ &
$2.53 (-3)$ \\
$40$ & $10$ & $1.66$ & $2.01\phantom{(-1)}$ & $1.02\phantom{(-1)}$ &
$2.42 (-2)$ & $8.98 (-3)$ \\
$40$ & $15$ & $2.45$ & $8.29\phantom{(-1)}$ & $2.81\phantom{(-1)}$ &
$4.02 (-1)$ & $1.30 (-1)$ \\
\hline
$80$ & $\phantom{1}1$  & $0.12$ & $4.10 (-1)$ & $4.07 (-1)$ & $2.87 (-5)$ &
$7.97 (-4)$ \\
$80$ & $\phantom{1}5$ & $0.61$ & $7.92 (-1)$ &  $7.19 (-1)$ & $4.57 (-4)$ &
$9.77 (-4)$ \\
$80$ & $10$ & $1.17$ & $1.22\phantom{(-1)}$ & $8.50 (-1)$ & $2.53 (-3)$ &
$1.48 (-3)$ \\
$80$ & $15$ & $1.73$ & $2.19\phantom{(-1)}$ & $1.06\phantom{(-1)}$ &
$1.52 (-2)$ & $5.51 (-3)$ \\
$80$ & $20$ & $2.29$ & $5.65\phantom{(-1)}$ & $2.04\phantom{(-1)}$ &
$1.07 (-1)$ & $3.53 (-2)$ \\
$80$ & $21$ & $2.41$ & $7.17\phantom{(-1)}$  & $2.48\phantom{(-1)}$ &
$1.60 (-1)$ & $5.21 (-2)$ \\
\hline
\end{tabular}
\end{center}
\end{table}

\begin{table}
\caption{Growth of $\phi^{(0)}(t):=t^\frac{1}{2} I_{0}(t^{2}/2)$}
\begin{center}
\begin{tabular}{|cc|}
\hline
 & \\
$t$ & $\phi^{(0)}(t)$\\
 & \\
\hline
$2.0$ & $3.22\phantom{(2)}$ \\
$2.5$ & $8.57\phantom{(2)}$ \\
$3.0$ & $3.03 (1)$ \\
$3.5$ & $1.41 (2)$ \\
$4.0$ & $8.55 (2)$ \\
\hline
\end{tabular}
\end{center}
\end{table}

\begin{table}
\caption{$m = 81$ and $\zeta=1$}
\begin{center}
\begin{tabular}{|cccc|}
\hline
& & & \\
$t$ & $\phi^{(1)}(t)$ & $n-1$ & $\phi_{n-1}(x_{m})/\sqrt{2m^{\frac{1}{2}}}$\\
& & & \\
\hline
$0.5$ & $6.24 (-1)$ & $\phantom{1}4$ & $6.25 (-1)$ \\
$1.5$ & $3.98 (-1)$ & $13$ & $3.96 (-1)$ \\
$2.5$ & $6.95 (-2)$ & $22$ & $6.68 (-2)$ \\
$3.5$ & $4.09 (-3)$ & $31$ & $3.48 (-3)$ \\
$4.5$ & $8.50 (-5)$ & $40$ & $5.32 (-5)$ \\
$5.5$ & $1.33 (-7)$ & $49$ & $2.07 (-7)$ \\
\hline
\end{tabular}
\end{center}
\end{table}

\end{document}







